%% 毕业设计小结

\begin{designsummary}

如果把毕业设计比作穿越一条未知的隧道，最初站在入口时，我看到的只有微光。

从选题开始就体会到"没有标准答案"的重量。课程作业有明确的输入输出和评判标准，而毕业设计从一开始就像在一张白纸上画地图——研究方向需要自己调研和锁定，技术路线需要在多个可能之间反复权衡。这个过程让我第一次感受到，真正的研究不是沿着既定的路走得更快，而是在不确定中找到值得走的路。

整个工作中最有收获的阶段，是将散落在不同论文中的方法串联为统一系统的过程。每篇论文都有自己假设的理想条件，但当它们被放在一起运行时，上游输出的精度波动、不同表征空间之间的语义鸿沟，都会产生论文中不曾讨论的摩擦。解决这些摩擦的过程让我体会到，做系统研究的核心能力不是单独优化某个模块，而是在模块之间的缝隙中找到让整体顺畅运转的衔接方案。

阅读论文的心得也在这次经历中发生了转变。以前读论文更关注方法本身是否有趣，现在会本能地去想这个方法对噪声的容忍度是多少、依赖的上游输出质量需要达到什么水平、和我正在用的其他模块能否兼容。这种带着具体问题阅读文献的方式，比泛泛浏览效率高得多，也更容易分辨出哪些技术细节是核心贡献、哪些是工程优化。

另一个感触是，科研中的"失败"和课程中的"做错"完全是两回事。课程作业做错了意味着扣分，而科研中很多看似没有产出的一两周，其实是在排除不可行的方向和积累对问题本质的理解。那些最终没有被采纳的方案，同样在帮助我逼近正确答案。这种对"进展"的重新定义，是我思维方式上重要的转变。

如今隧道走到了尽头，那些熬夜推敲、反复试错的日子，最终化作了这段文字里平静的回望。这段经历带给我的，远不只是完成了某个系统或技术方案，更是一种敢于面对未知、在漫长反馈周期中持续前行的底气。

\end{designsummary}